\subsection{比較判定法}

$0\leq f(x)\leq g(x)$が端点の十分近くで成り立つとする.このとき,\,$g$の広義積分が収束すれば$f$の広義積分も収束する.逆向きには,\,$f$の広義積分が$+\infty$へ発散すれば,\,$g$の広義積分も$+\infty$へ発散する.

例えば無限区間では,\,$\int_A^R f(x)\dd x$は$R$とともに増加する.しかも
\[
 0\leq\int_A^Rf(x)\dd x\leq\int_A^Rg(x)\dd x
 \leq\int_A^{\infty}g(x)\dd x
\]
と一定の上限をもつ.上に有界な単調増加の量は有限の極限をもつので,\,$f$の積分も収束する.端点付近でだけ比較できればよいのは,それ以外の閉区間の積分は有限の一定値を加えるだけだからである.

特に$f,g>0$で
\[
 \lim\frac{f(x)}{g(x)}=c,\qquad0<c<\infty
\]
となるなら,端点の十分近くで$\frac c2g(x)\leq f(x)\leq\frac{3c}{2}g(x)$と挟める.そのため二つの積分は同時に収束するか,同時に発散する.この極限比較では,\,$c$が正の有限値であるという条件が大切である.比の極限が$0$の場合は,\,$f$が$g$より小さいことしか分からず,両者の収束が一致するとは限らない.

\begin{example}{広義積分の置換}
$\int_0^1\frac1{\sqrt{x(1+x)}}\dd x$の収束を確かめ,値を求めよ.
\end{example}

特異性があるのは$x=0$だけである.\,$0<x\leq1$で
\[
 0<\frac1{\sqrt{x(1+x)}}\leq\frac1{\sqrt{x}}
\]
となり,右辺は$0$の近くで積分可能だから,求める広義積分は収束する.

計算では,まず$\varepsilon\leq x\leq1$に限って$x=t^2$と置く.平方根の$x$と$\dd x=2t\dd t$が約分されるので
\[
 \int_{\varepsilon}^1\frac1{\sqrt{x(1+x)}}\dd x
 =2\int_{\sqrt{\varepsilon}}^1\frac1{\sqrt{1+t^2}}\dd t.
\]
右辺の被積分関数は$t=0$でも連続である.\,$\varepsilon\to0+$として
\[
 2\int_0^1\frac1{\sqrt{1+t^2}}\dd t
 =2\left[\log\left(t+\sqrt{1+t^2}\right)\right]_0^1
 =2\log(1+\sqrt2)
\]
を得る.置換はまず通常の定積分に行い,その後で端点の極限をとった.


